\pagestyle{empty}
\tight
\begin{tblr}{width=\textwidth,colspec={|Q[c,m,1.9cm]|X[l,m]|},hlines,vlines,hspan=minimal,row{1}={bg=headblue},rowsep=1pt,colsep=3pt}
\SetCell{c,m}\hfil\logo\hfil & \textbf{POLITEKNIK NEGERI MALANG}\par\textbf{JURUSAN TEKNOLOGI INFORMASI}\par\textbf{PROGRAM STUDI: D4 TEKNIK INFORMATIKA} \\
\SetCell[c=2]{c} \textbf{RENCANA PEMBELAJARAN SEMESTER (RPS)} & \\
\end{tblr}
\begin{longtblr}{width=\textwidth,colspec={|Q[l,m,1.9cm]|X[0.85,l,m]|X[1.45,l,m]|X[0.75,l,m]|X[1.15,l,m]|},hlines,vlines,hspan=minimal,rowsep=1pt,colsep=2pt,row{1,3}={bg=headgray,font=\bfseries}}
MATA KULIAH & KODE & BOBOT (SKS)/jam & SEMESTER & TGL. PENYUSUNAN \\
Pemrograman Berbasis Framework & RTI256003 & 2 SKS/4 jam & 6 & 15 April 2026 \\
OTORISASI & Dosen Pengembang RPS & Koordinator RMK & \SetCell[c=2]{l} Ka PRODI &  \\
 & Farid Angga Pribadi, S.Kom., M.Kom. &  & \SetCell[c=2]{l}{\vspace{8mm}\par Dr. Ely Setyo Astuti, S.T., M.T.} &  \\
\SetCell[r=10]{l} Capaian Pembelajaran (CP) & \SetCell[c=4]{l,bg=headgray,font=\bfseries} Capaian Pembelajaran Lulusan Program Studi (CPL-Prodi) &  &  &  \\
 & CPL02 & \SetCell[c=3]{l} Mampu menerapkan prinsip rekayasa sistem dan perangkat lunak untuk menghasilkan solusi aplikasi multi-platform sesuai kebutuhan. &  &  \\
 & CPL06 & \SetCell[c=3]{l} Mampu merancang, mengimplementasikan, menguji, melakukan deployment, memelihara, serta menjamin mutu perangkat lunak yang menjawab permasalahan dan memenuhi kebutuhan pengguna. &  &  \\
 & \SetCell[c=4]{l,bg=headgray,font=\bfseries} Capaian Pembelajaran mata kuliah (CPL-MK) &  &  &  \\
 & CPMK0209 & \SetCell[c=3]{l} Mahasiswa mampu memilih dan menerapkan teknologi pengembangan berbasis platform untuk membangun aplikasi sesuai kebutuhan. &  &  \\
 & CPMK0603 & \SetCell[c=3]{l} Mahasiswa mampu mengimplementasikan dan melakukan deployment aplikasi web, mobile, atau framework sesuai kebutuhan pengguna. &  &  \\
 & \SetCell[c=4]{l,bg=headgray,font=\bfseries} Kemampuan akhir tiap tahapan belajar (Sub-CPMK) &  &  &  \\
 & SCPMK0209-04701 & \SetCell[c=3]{l} Mahasiswa mampu menerapkan setup, routing, styling, custom document, API routes, CSR, SSR, SSG, ISR, middleware, dan route protection. &  &  \\
 & SCPMK0603-04702 & \SetCell[c=3]{l} Mahasiswa mampu mengimplementasikan authentication, multi-role system, optimasi, dan unit testing pada aplikasi framework modern. &  &  \\
 & SCPMK0603-04703 & \SetCell[c=3]{l} Mahasiswa mampu menghasilkan aplikasi berbasis framework melalui PBL berdasarkan kebutuhan stakeholder. &  &  \\
\SetCell[r=2]{l} Penilaian & \SetCell[c=4]{l} *Nilai CPMK didapat dari agregasi nilai SUB-CPMK yang dibebankan &  &  &  \\
 & \SetCell[c=4]{l}{\tiny\begin{tblr}{width=\linewidth,colspec={|Q[c,m,0.1\linewidth]|Q[c,m,0.16\linewidth]|X[l,m]|X[l,m]|X[l,m]|Q[c,m,0.08\linewidth]|},hlines,vlines,hspan=minimal,rowsep=1pt,colsep=0.7pt,row{1,2}={bg=headgray,font=\bfseries}}
Kode CPMK & Kode SCPMK & \SetCell[c=3]{c} Bobot per Bentuk Penilaian & & & Total Bobot \\
 & & Tugas & UTS & PBL & \\
\SetCell[r=1]{c} CPMK0209 & SCPMK0209-04701 & 10\%: setup, routing, API, SSR/SSG/ISR & 10\%: framework fundamentals & & 20.00\% \\
\SetCell[r=2]{c} CPMK0603 & SCPMK0603-04702 & 10\%: auth, multi-role, optimasi, unit testing & & & 10.00\% \\
 & SCPMK0603-04703 & & & 70\%: PBL aplikasi berbasis framework & 70.00\% \\
\SetCell[c=5]{r}\textbf{Total Per Penilaian} & & & & & \textbf{100\%} \\
\end{tblr}} &  &  &  \\
Diskripsi Singkat Mata Kuliah & \SetCell[c=4]{l} Mata kuliah Pemrograman Berbasis Framework membahas pengembangan aplikasi web modern berbasis framework JavaScript/TypeScript, mencakup setup, routing, styling, API routes, CSR, SSR, SSG, ISR, middleware, authentication, multi-role, optimasi, unit testing, dan proyek PBL. &  &  &  \\
Bahan Kajian / Materi Pembelajaran & \SetCell[c=4]{l}\begin{enumerate}\item Setup, routing, styling, custom document\item API routes, CSR, SSR, SSG, ISR, middleware, route protection\item Authentication, multi-role, optimasi, unit testing, dan proyek PBL\end{enumerate} &  &  &  \\
\SetCell[r=2]{l} Pustaka & \SetCell[c=4]{l} \textbf{Utama:}\begin{enumerate}\item Next.js Documentation. https://nextjs.org/docs.\item React Documentation. https://react.dev.\item Testing Library Documentation. https://testing-library.com.\end{enumerate} &  &  &  \\
 & \SetCell[c=4]{l} \textbf{Pendukung:} Materi LMS, dokumentasi resmi perangkat lunak, dan jobsheet praktikum. &  &  &  \\
Media Pembelajaran & \SetCell[c=2]{l} \textbf{Software:} Node.js, npm/pnpm, framework web modern, Git, VS Code, browser, testing tools. &  & \SetCell[c=2]{l} \textbf{Hardware:} Komputer/laptop dan proyektor. &  \\
Nama Dosen Pengampu & \SetCell[c=4]{l}\begin{enumerate}\item Farid Angga Pribadi, S.Kom., M.Kom.\item Dimas Wahyu Wibowo, S.T., M.T.\item Muhammad Unggul Pamenang, S.St., M.T.\item Dian Hanifudin Subhi, S.Kom., M.Kom.\end{enumerate} &  &  &  \\
Matakuliah Syarat & \SetCell[c=4]{l} Pemrograman Web Lanjut &  &  &  \\
\end{longtblr}
\begin{longtblr}{width=\textwidth,colspec={|Q[c,m,0.06\textwidth]|X[1.35,l,m]|X[1.08,l,m]|X[1.16,l,m]|X[0.7,l,m]|X[1.24,l,m]|X[1.06,l,m]|X[0.98,l,m]|Q[c,m,0.05\textwidth]|},hlines,vlines,hspan=minimal,rowhead=2,row{1,2}={bg=headgreen,font=\bfseries},rowsep=1pt,colsep=1.5pt}
Mgg.\par Ke & Kemampuan Akhir Yang Direncanakan (Sub-CP-MK) & Bahan kajian (Materi Pembelajaran) & Bentuk dan Metode Pembelajaran & Estimasi Waktu & Pengalaman Belajar Mahasiswa & Kriteria \& Bentuk Penilaian & Indikator Penilaian & Bobot Penilaian (\%) \\
(1) & (2) & (3) & (4) & (5) & (6) & (7) & (8) & (9) \\
1 & SCPMK0209-04701 & Setup dan routing. & Modalitas: Blended Learning. Bentuk: Praktikum/PBL. Metode: hands-on practice, mentoring, code review, demo, dan presentasi. & 1 x 2 x 50' tatap muka; 1 x 2 x 50' tugas/praktik mandiri. & Mahasiswa mengerjakan aktivitas terarah untuk setup dan routing dan menunjukkan evidence hasil belajar. & Bentuk penilaian: Setup dan routing. Mengacu rubrik RTM. & Ketepatan konsep; kualitas artefak/praktik; validasi dan dokumentasi. & 2 \\
2 & SCPMK0209-04701 & Styling dan custom document. & Modalitas: Blended Learning. Bentuk: Praktikum/PBL. Metode: hands-on practice, mentoring, code review, demo, dan presentasi. & 1 x 2 x 50' tatap muka; 1 x 2 x 50' tugas/praktik mandiri. & Mahasiswa mengerjakan aktivitas terarah untuk styling dan custom document dan menunjukkan evidence hasil belajar. & Bentuk penilaian: Styling dan custom document. Mengacu rubrik RTM. & Ketepatan konsep; kualitas artefak/praktik; validasi dan dokumentasi. & 2 \\
3 & SCPMK0209-04701 & API routes dan CSR. & Modalitas: Blended Learning. Bentuk: Praktikum/PBL. Metode: hands-on practice, mentoring, code review, demo, dan presentasi. & 1 x 2 x 50' tatap muka; 1 x 2 x 50' tugas/praktik mandiri. & Mahasiswa mengerjakan aktivitas terarah untuk api routes dan csr dan menunjukkan evidence hasil belajar. & Bentuk penilaian: API routes dan CSR. Mengacu rubrik RTM. & Ketepatan konsep; kualitas artefak/praktik; validasi dan dokumentasi. & 2 \\
4 & SCPMK0209-04701 & SSR dan static generation. & Modalitas: Blended Learning. Bentuk: Praktikum/PBL. Metode: hands-on practice, mentoring, code review, demo, dan presentasi. & 1 x 2 x 50' tatap muka; 1 x 2 x 50' tugas/praktik mandiri. & Mahasiswa mengerjakan aktivitas terarah untuk ssr dan static generation dan menunjukkan evidence hasil belajar. & Bentuk penilaian: SSR dan static generation. Mengacu rubrik RTM. & Ketepatan konsep; kualitas artefak/praktik; validasi dan dokumentasi. & 2 \\
5 & SCPMK0209-04701 & Dynamic routing dan ISR. & Modalitas: Blended Learning. Bentuk: Praktikum/PBL. Metode: hands-on practice, mentoring, code review, demo, dan presentasi. & 1 x 2 x 50' tatap muka; 1 x 2 x 50' tugas/praktik mandiri. & Mahasiswa mengerjakan aktivitas terarah untuk dynamic routing dan isr dan menunjukkan evidence hasil belajar. & Bentuk penilaian: Dynamic routing dan ISR. Mengacu rubrik RTM. & Ketepatan konsep; kualitas artefak/praktik; validasi dan dokumentasi. & 2 \\
6 & SCPMK0209-04701 & Middleware dan route protection. & Modalitas: Blended Learning. Bentuk: Praktikum/PBL. Metode: hands-on practice, mentoring, code review, demo, dan presentasi. & 1 x 2 x 50' tatap muka; 1 x 2 x 50' tugas/praktik mandiri. & Mahasiswa mengerjakan aktivitas terarah untuk middleware dan route protection dan menunjukkan evidence hasil belajar. & Bentuk penilaian: Middleware dan route protection. Mengacu rubrik RTM. & Ketepatan konsep; kualitas artefak/praktik; validasi dan dokumentasi. & 0 \\
7 & SCPMK0603-04702 & Authentication dan multi role. & Modalitas: Blended Learning. Bentuk: Praktikum/PBL. Metode: hands-on practice, mentoring, code review, demo, dan presentasi. & 1 x 2 x 50' tatap muka; 1 x 2 x 50' tugas/praktik mandiri. & Mahasiswa mengerjakan aktivitas terarah untuk authentication dan multi role dan menunjukkan evidence hasil belajar. & Bentuk penilaian: Authentication dan multi role. Mengacu rubrik RTM. & Ketepatan konsep; kualitas artefak/praktik; validasi dan dokumentasi. & 0 \\
8 & SCPMK0603-04702 & Optimization dan unit testing. & Modalitas: Blended Learning. Bentuk: Praktikum/PBL. Metode: hands-on practice, mentoring, code review, demo, dan presentasi. & 1 x 2 x 50' tatap muka; 1 x 2 x 50' tugas/praktik mandiri. & Mahasiswa mengerjakan aktivitas terarah untuk optimization dan unit testing dan menunjukkan evidence hasil belajar. & Bentuk penilaian: Optimization dan unit testing. Mengacu rubrik RTM. & Ketepatan konsep; kualitas artefak/praktik; validasi dan dokumentasi. & 0 \\
9 & SCPMK0209-04701, SCPMK0603-04702 & UTS framework fundamentals. & Modalitas: Blended Learning. Bentuk: Praktikum/PBL. Metode: hands-on practice, mentoring, code review, demo, dan presentasi. & 1 x 2 x 50' tatap muka; 1 x 2 x 50' tugas/praktik mandiri. & Mahasiswa mengerjakan aktivitas terarah untuk uts framework fundamentals dan menunjukkan evidence hasil belajar. & Bentuk penilaian: UTS framework fundamentals. Mengacu rubrik RTM. & Ketepatan konsep; kualitas artefak/praktik; validasi dan dokumentasi. & 20 \\
10 & SCPMK0603-04703 & PBL analisis kebutuhan. & Modalitas: Blended Learning. Bentuk: Praktikum/PBL. Metode: hands-on practice, mentoring, code review, demo, dan presentasi. & 1 x 2 x 50' tatap muka; 1 x 2 x 50' tugas/praktik mandiri. & Mahasiswa mengerjakan aktivitas terarah untuk pbl analisis kebutuhan dan menunjukkan evidence hasil belajar. & Bentuk penilaian: PBL analisis kebutuhan. Mengacu rubrik RTM. & Ketepatan konsep; kualitas artefak/praktik; validasi dan dokumentasi. & 10 \\
11 & SCPMK0603-04703 & PBL fitur inti. & Modalitas: Blended Learning. Bentuk: Praktikum/PBL. Metode: hands-on practice, mentoring, code review, demo, dan presentasi. & 1 x 2 x 50' tatap muka; 1 x 2 x 50' tugas/praktik mandiri. & Mahasiswa mengerjakan aktivitas terarah untuk pbl fitur inti dan menunjukkan evidence hasil belajar. & Bentuk penilaian: PBL fitur inti. Mengacu rubrik RTM. & Ketepatan konsep; kualitas artefak/praktik; validasi dan dokumentasi. & 10 \\
12 & SCPMK0603-04703 & PBL data/API dan state. & Modalitas: Blended Learning. Bentuk: Praktikum/PBL. Metode: hands-on practice, mentoring, code review, demo, dan presentasi. & 1 x 2 x 50' tatap muka; 1 x 2 x 50' tugas/praktik mandiri. & Mahasiswa mengerjakan aktivitas terarah untuk pbl data/api dan state dan menunjukkan evidence hasil belajar. & Bentuk penilaian: PBL data/API dan state. Mengacu rubrik RTM. & Ketepatan konsep; kualitas artefak/praktik; validasi dan dokumentasi. & 10 \\
13 & SCPMK0603-04703 & PBL auth/role. & Modalitas: Blended Learning. Bentuk: Praktikum/PBL. Metode: hands-on practice, mentoring, code review, demo, dan presentasi. & 1 x 2 x 50' tatap muka; 1 x 2 x 50' tugas/praktik mandiri. & Mahasiswa mengerjakan aktivitas terarah untuk pbl auth/role dan menunjukkan evidence hasil belajar. & Bentuk penilaian: PBL auth/role. Mengacu rubrik RTM. & Ketepatan konsep; kualitas artefak/praktik; validasi dan dokumentasi. & 10 \\
14 & SCPMK0603-04703 & PBL testing dan optimasi. & Modalitas: Blended Learning. Bentuk: Praktikum/PBL. Metode: hands-on practice, mentoring, code review, demo, dan presentasi. & 1 x 2 x 50' tatap muka; 1 x 2 x 50' tugas/praktik mandiri. & Mahasiswa mengerjakan aktivitas terarah untuk pbl testing dan optimasi dan menunjukkan evidence hasil belajar. & Bentuk penilaian: PBL testing dan optimasi. Mengacu rubrik RTM. & Ketepatan konsep; kualitas artefak/praktik; validasi dan dokumentasi. & 10 \\
15 & SCPMK0603-04703 & PBL deployment preparation. & Modalitas: Blended Learning. Bentuk: Praktikum/PBL. Metode: hands-on practice, mentoring, code review, demo, dan presentasi. & 1 x 2 x 50' tatap muka; 1 x 2 x 50' tugas/praktik mandiri. & Mahasiswa mengerjakan aktivitas terarah untuk pbl deployment preparation dan menunjukkan evidence hasil belajar. & Bentuk penilaian: PBL deployment preparation. Mengacu rubrik RTM. & Ketepatan konsep; kualitas artefak/praktik; validasi dan dokumentasi. & 10 \\
16 & SCPMK0603-04703 & PBL finalisasi proyek. & Modalitas: Blended Learning. Bentuk: Praktikum/PBL. Metode: hands-on practice, mentoring, code review, demo, dan presentasi. & 1 x 2 x 50' tatap muka; 1 x 2 x 50' tugas/praktik mandiri. & Mahasiswa mengerjakan aktivitas terarah untuk pbl finalisasi proyek dan menunjukkan evidence hasil belajar. & Bentuk penilaian: PBL finalisasi proyek. Mengacu rubrik RTM. & Ketepatan konsep; kualitas artefak/praktik; validasi dan dokumentasi. & 10 \\
17 & SCPMK0603-04703 & UAS presentasi hasil proyek. & Modalitas: Blended Learning. Bentuk: Praktikum/PBL. Metode: hands-on practice, mentoring, code review, demo, dan presentasi. & 1 x 2 x 50' tatap muka; 1 x 2 x 50' tugas/praktik mandiri. & Mahasiswa mengerjakan aktivitas terarah untuk uas presentasi hasil proyek dan menunjukkan evidence hasil belajar. & Bentuk penilaian: UAS presentasi hasil proyek. Mengacu rubrik RTM. & Ketepatan konsep; kualitas artefak/praktik; validasi dan dokumentasi. & 0 \\
\end{longtblr}
